\section{Theoretical Framework and Implemented Objective}

To systematically find adversarial prefixes, we rely on gradient-based
optimization over continuous token embeddings. However, optimizing a continuous
surrogate does not guarantee that its nearest discrete projection remains
adversarial. In this section, we formalize the optimization objective,
acknowledge the theoretical limits of unconditional projection, and establish
the precise margin conditions under which our continuous-to-discrete pipeline is
guaranteed to succeed.

\subsection{Setup and the Continuous Attack Space}

Let the attacker optimize a prefix of length $m$. A discrete prefix is a
sequence $z = (z_1, \dots, z_m) \in \mathcal{V}^m$, where $\mathcal{V}$ is the
vocabulary. Let $E \in \mathbb{R}^{|\mathcal{V}| \times d}$ denote the
embedding matrix. We define a continuous relaxation
$U = (u_1, \dots, u_m)$, where each $u_t \in \Delta^{|\mathcal{V}|}$ is a
probability simplex vector. The relaxed embedding at position $t$ is
\[
e_t(U) = E^\top u_t.
\]

Let $p_i^U$ and $q_i^U$ denote the target and draft next-token distributions at
speculative step $i$ under the relaxed prefix $U$. Drawing from the one-step
expected acceptance identity $\mathbb{E}[\alpha_i] = 1 - \tv(p_i, q_i)$, a
natural theoretical objective is to maximize early draft-target disagreement:
\[
C(U) = \sum_{i=1}^{\gamma} w_i \tv(p_i^U, q_i^U),
\]
where $w_i = \gamma - i + 1$ linearly weights earlier tokens to maximize
rejection pressure. The corresponding discrete objective is defined by
restriction:
\[
C_d(z) := C(E(z)),
\]
where $E(z)$ is the concatenated embedding sequence of the discrete tokens.

To recover a deployable attack, we define the token-wise nearest-neighbor
projection $\Pi(U) := (\Pi_1(U), \dots, \Pi_m(U))$, where
\[
\Pi_t(U) = \arg\min_{v \in \mathcal{V}} \|e_t(U) - E_v\|_2.
\]
We quantify the continuous-to-discrete shift via the projection distortion
$\rho(U)$:
\[
\rho(U) := \sum_{t=1}^{m} \|e_t(U) - E_{\Pi_t(U)}\|_2.
\]

\subsection{The Obstruction of Prompt Waywardness}

A critical vulnerability in gradient-based text attacks is the assumption that
continuous optimization smoothly maps back to discrete space. We first state a
negative result demonstrating that without structural assumptions, projection is
fundamentally unsafe.

\begin{proposition}[Prompt Waywardness]
\label{prop:waywardness}
For any finite token embedding space with at least two tokens and any
$\delta > 0$, there exists a smooth objective $C$ over the convex hull of token
embeddings and a relaxed optimizer $U^*$ such that $U^*$ is within distance
$\delta$ of a discrete token cell, yet the nearest-neighbor projection
$\Pi(U^*)$ is arbitrarily suboptimal among all discrete prefixes.
\end{proposition}

This proposition formalizes the empirical warning signs from prior literature:
an unconditional projection guarantee is impossible. To bind the discrete gap,
we must constrain the behavior of the specific model distributions.

\subsection{Bounding the Projection Gap}

To guarantee attack transferability from continuous to discrete space, we assume
the generative distributions exhibit Lipschitz continuity with respect to the
input embeddings.

\begin{assumption}[TV-Lipschitzness]
\label{assump:tv-lipschitz}
For each speculative step $i$, there exist constants $L_i^p, L_i^q$ such that
for any two relaxed prefixes $U, V$,
\[
\tv(p_i^U, p_i^V) \le L_i^p \|U - V\|_{2,1},
\]
\[
\tv(q_i^U, q_i^V) \le L_i^q \|U - V\|_{2,1},
\]
where
\[
\|U - V\|_{2,1} := \sum_{t=1}^{m} \|e_t(U) - e_t(V)\|_2.
\]
Let
\[
L_C := \sum_{i=1}^{\gamma} w_i (L_i^p + L_i^q)
\]
denote the aggregate Lipschitz constant of the weighted objective.
\end{assumption}

Under this assumption, we can strictly bound the quality degradation incurred
during projection.

\begin{theorem}[Relaxation Gap for Early Acceptance Collapse]
\label{thm:relax-gap}
Let $U^*$ be an $\varepsilon$-optimal solution of the relaxed problem, such
that
\[
C(U^*) \ge \sup_U C(U) - \varepsilon.
\]
Let
\[
z^* \in \arg\max_{z \in \mathcal{V}^m} C_d(z)
\]
be the optimal discrete prefix. Then the projected discrete attack satisfies
\[
C_d(\Pi(U^*)) \ge C_d(z^*) - \varepsilon - L_C \rho(U^*).
\]
\end{theorem}

\begin{proof}
By Assumption~\ref{assump:tv-lipschitz} and the triangle inequality, $C$ is
$L_C$-Lipschitz. Let
$\hat{z} = \Pi(U^*)$. Because the continuous embeddings of $\hat{z}$ are within
$\rho(U^*)$ of $U^*$, it follows that
\[
C_d(\hat{z}) = C(E(\hat{z})) \ge C(U^*) - L_C \rho(U^*).
\]
Since the optimal discrete prompt is a valid feasible solution in the relaxed
space,
\[
C(U^*) \ge C_d(z^*) - \varepsilon.
\]
Combining these inequalities yields the bound.
\end{proof}

Theorem~\ref{thm:relax-gap} establishes that the discrete attack gap is
controlled by only two
actionable quantities: the optimization error of the relaxed solver
($\varepsilon$) and the distance from the relaxed solution to the token
manifold ($\rho(U^*)$).

We can further characterize when the continuous relaxation exactly recovers the
globally optimal discrete attack. Let the discrete attack margin be defined as
\[
\Gamma := C_d(z^*) - \max_{z \neq z^*} C_d(z).
\]

\begin{corollary}[Exact Recovery under Discrete Margin]
\label{cor:exact-recovery}
Assume the optimal discrete adversarial prefix $z^*$ is unique
($\Gamma > 0$). Under Theorem~\ref{thm:relax-gap}, if the combined optimization
error and
projection distortion is strictly less than the discrete margin,
\[
\varepsilon + L_C \rho(U^*) < \Gamma,
\]
then the nearest-neighbor projection perfectly identifies the optimal discrete
attack:
\[
\Pi(U^*) = z^*.
\]
\end{corollary}

\subsection{Algorithm Design Justification}

The theoretical bounds above directly dictate our optimization strategy. To
maximize the discrete adversarial effect $C_d(\Pi(U))$, we cannot simply
maximize $C(U)$. We must explicitly minimize the projection gap
$L_C \rho(U)$. Consequently, any practical attack should optimize a penalized
continuous objective of the form
\[
\max_U \; C(U) - \lambda \rho(U) - \tau \sum_{t=1}^{m} H(u_t),
\]
where $\lambda \rho(U)$ acts as a manifold-attraction penalty to explicitly
bound the gap in Theorem~\ref{thm:relax-gap}, and the entropy regularization
sharpens the simplex vectors so the optimizer explores regions near discrete
token boundaries.

\subsection{Implemented Soft-Prefix Attack}

Our implementation uses the same relaxed search space, but instantiates the
continuous objective for the \emph{first verifier step}, because early rejection
creates the largest downstream slowdown in token-wise speculative decoding. Let
$p^U$ and $q^U$ denote the target and draft next-token distributions after the
relaxed prefix $U$. We define the implemented objective
\begin{align*}
J(U) ={}&
\underbrace{\sum_{y} q^U(y)\,\mathrm{softplus}\!\left(\log q^U(y)-\log p^U(y)\right)}_{\text{soft collapse}} \\
&+ \beta \,\tv(p^U, q^U)
+ \eta \,\mathrm{KL}(q^U \Vert p^U) \\
&- \tau \,\frac{1}{m}\sum_{t=1}^{m} H(u_t)
- \lambda \,\rho(U).
\end{align*}
The first term directly rewards the failure mode we seek: tokens on which the
draft is confident while the verifier assigns much lower probability. The TV
and reverse-KL terms enlarge the broader disagreement region around that token,
and the entropy term sharpens the relaxed simplex vectors so that projection is
stable. The projection penalty $\rho(U)$ is the theory-guided term from
Theorem~\ref{thm:relax-gap}.

We optimize $J(U)$ with Adam over per-position vocabulary logits. Each step
uses a straight-through argmax estimator: the forward pass uses the projected
hard token at each position, while the backward pass differentiates through the
softmax relaxation. This is important operationally. The optimizer therefore
sees the same discrete token manifold that deployment will use, instead of
optimizing a purely continuous embedding string that may not survive
projection.

While Theorem~\ref{thm:relax-gap} bounds the discrete gap using explicit
$L_2$ nearest-neighbor projection in the embedding space, evaluating pairwise
Euclidean distances across the full tokenizer vocabulary at every optimization
step is computationally prohibitive. In the implementation, we therefore
approximate the operational projection with a straight-through argmax estimator
over the relaxed probability simplex. The theoretical connection is preserved
by retaining the explicit $L_2$ projection-distortion penalty $\rho(U)$ in the
continuous objective, which penalizes the distance between the relaxed
soft-embeddings and the hard embeddings induced by the argmax-selected tokens.

Model selection is also verifier-aligned. Although optimization uses the smooth
objective $J(U)$, we keep the best candidate by its exact projected
first-token acceptance
\[
\alpha_1 = \min\!\left(1, \frac{p(y_1)}{q(y_1)}\right),
\]
computed after hard projection. In other words, the smooth objective drives the
search, but the retained attack is always evaluated in the exact discrete
speculative-decoding pipeline.

In the current experiments, the key hyperparameters are a five-token prefix,
Adam learning rate $0.3$, temperature $1.0$, entropy weight $0.01$, TV weight
$1.0$, reverse-KL weight $0.1$, and gradient clipping at $1.0$. The
projection-distortion weight $\lambda$ is calibrated automatically from the
initial ratio between the unpenalized objective and the initial projection
distortion, then swept on Palmetto through a multiplicative factor. This
calibration gives a stable way to compare attacks across different prefix
lengths without hand-retuning the objective scale each time. Because the token
embedding space is highly non-convex, initializing the relaxed prefix from a
uniform distribution is unstable. We therefore warm-start the optimizer's
initial logits from a manual adversarial seed and then search the full
continuous vocabulary space around that initialization. In practice, the search
often discovers subword fragments or visually unusual tokens that strictly
outperform the original seed.
